\documentclass[12pt,a4paper]{article}
\usepackage[utf8]{inputenc}
\usepackage[french]{babel}
\usepackage[T1]{fontenc}
\usepackage{geometry}
\geometry{margin=2.5cm}
\usepackage{titlesec}
\usepackage{enumitem}
\usepackage{xcolor}
\usepackage{fancyhdr}

% Configuration des titres
\titleformat{\section}{\LARGE\textbf{\color{blue!80!black}}\centering}{}{0em}{}[\vspace{1ex}\titlerule]
\titleformat{\subsection}{\Large\textbf{\color{blue!60!black}}}{}{0em}{\medskip}
\titleformat{\subsubsection}{\large\textbf}{}{0em}{\smallskip}

% Mise en page
\pagestyle{fancy}
\fancyhf{}
\lhead{Obscurus Anti-Cheat Solutions}
\rhead{Fiche de Poste H/F - 2026}
\cfoot{\thepage}

\begin{document}

\section{RESPONSABLE SÉCURITÉ OFFENSIVE (H/F)}

\subsection{Identification du Poste}
\begin{itemize}[label=\textbullet]
    \item \textbf{Département} : Technique / Sécurité
    \item \textbf{Supérieur} : Directeur Général \& CTO
    \item \textbf{Collaborations} : Équipe commerciale, RH
    \item \textbf{Rôle} : Tester les défenses du logiciel en jouant le rôle de l'attaquant
\end{itemize}

\subsection{Missions Principales}

\subsubsection{Tests d'Attaque}
\begin{itemize}[label=\textbullet]
    \item Créer des outils de triche en interne (aimbot, wallhack) pour tester les détections.
    \item Tenter de contourner le logiciel anti-cheat avec des techniques avancées.
    \item Tester la résistance du système face à des données falsifiées.
\end{itemize}

\subsubsection{Analyse et Veille}
\begin{itemize}[label=\textbullet]
    \item Étudier les cheats disponibles sur le marché pour comprendre leurs méthodes.
    \item Auditer le code de l'entreprise pour repérer des failles potentielles.
\end{itemize}

\subsubsection{Collaboration avec les Développeurs}
\begin{itemize}[label=\textbullet]
    \item Rédiger des rapports clairs sur les failles trouvées, avec niveau de dangerosité.
    \item Travailler avec l'équipe de développement pour corriger et renforcer le logiciel.
\end{itemize}

\subsection{Compétences Requises}
\begin{itemize}[label=\textbullet]
    \item \textbf{Technique} : Maîtrise des outils d'analyse (IDA Pro, Ghidra), développement en C++ et Python.
    \item \textbf{Humaines} : Curiosité, éthique rigoureuse, esprit critique.
\end{itemize}

\subsection{Indicateurs de Réussite}
\begin{itemize}[label=\textbullet]
    \item Nombre de failles trouvées avant qu'elles ne soient exploitées à l'extérieur.
    \item Rapidité à reproduire un contournement signalé.
\end{itemize}

\end{document}
